\section{Marco Ético Aplicable}

Astro-Denoise se sitúa en la intersección de tres marcos normativos: la legislación colombiana de protección de datos, las políticas nacionales de ciencia abierta e inteligencia artificial, y los estándares internacionales de gobernanza de datos científicos.

\subsection{Normativa Colombiana}

La Ley 1581 de 2012 (\cite{colombia2012}) establece principios de transparencia, finalidad y veracidad que, aunque pensados para datos personales, fijan el estándar ético nacional para cualquier manejo de información. Astro-Denoise los cumple operando con datos simulados de acceso público, licencias explícitas y metadatos completos.

La Ley 2502 de 2025 (\cite{ley2502}) introduce regulación específica para sistemas de inteligencia artificial en Colombia, incluyendo requisitos de transparencia algorítmica y rendición de cuentas. El CONPES 4144 (\cite{conpes4144}) complementa esta normativa con una política nacional de IA que promueve investigación ética, desarrollo responsable y acceso equitativo. Los lineamientos de seguridad y privacidad del MinTIC (\cite{mintic2024}) establecen buenas prácticas para sistemas de IA que Astro-Denoise adopta en su diseño.

La Política de Ciencia Abierta de Minciencias (\cite{minciencias2022}) exige acceso abierto a resultados financiados con fondos públicos, disponibilidad permanente de datos y transparencia metodológica. Astro-Denoise encarna estos principios: código MIT, datos CC0, documentación pública y contribuciones abiertas.

\subsection{Estándares Internacionales}

Los principios FAIR (\cite{wilkinson2016}) — Findable, Accessible, Interoperable, Reusable — son el referente global para gobernanza de datos científicos. En astronomía, la alineación con los estándares del IVOA (International Virtual Observatory Alliance) hace que la implementación FAIR sea natural (\cite{otoole2022}). Astro-Denoise cumple cada principio: publicación con DOI, acceso sin autenticación, formatos JSON/CSV estándar, y licencias explícitas (\cite{schindler2024}).

La Recomendación de Ciencia Abierta de la UNESCO (\cite{unesco2021}) consagra el acceso abierto como derecho humano. El proyecto responde a este mandato mediante código y datos públicos, documentación en español, y una arquitectura que elimina barreras geográficas e institucionales (\cite{harrison2017}).

\subsection{Síntesis}

El marco ético de Astro-Denoise opera en tres niveles: nacional (Ley 1581, Ley 2502, CONPES 4144, Minciencias), regional (compromisos latinoamericanos de inclusión científica) y global (FAIR, UNESCO). Todos convergen en un mismo principio: la ciencia abierta como mecanismo de equidad y verificación.
